%! Author = k
%! Date = 2024/11/20

% Preamble
\documentclass[11pt]{article}

\usepackage{amsmath}
\usepackage[UTF8]{ctex}
\usepackage{multirow}
\usepackage{booktabs} % 用于更好看的表格线
\usepackage{array} % 用于调整列宽
\usepackage{caption} % 用于设置表格标题
\usepackage{geometry}

\pagestyle{empty} % 隐藏页眉页脚

\captionsetup[table]{labelformat=empty} % 禁用表格编号

\geometry{a4paper,paperwidth = 30cm, paperheight = 30cm}

\begin{document}

    % 表格代码开始
    \begin{table}[h]
        \centering
        \renewcommand{\arraystretch}{1.5} % 调整行间距
        \setlength{\tabcolsep}{12pt} % 调整列间距
        \begin{tabular}{>{\centering}p{2cm}|>{\centering}p{2cm}|>{\centering}p{2cm}|>{\centering}p{2cm}|c} % 控制边框：| 表示有竖线，c 表示居中对齐
            \hline % 表格顶部边框
            \textbf{$J$} & \textbf{$K$} & \textbf{$Q^{n}$} & \textbf{$Q^{n + 1}$} & \textbf{输出状态说明} \\ % 表头
            \hline % 表头下边框
            0 & 0 & 0 & \multirow{2}{*}{ $ \hspace{0.7cm} \left. \begin{aligned} 0 \\ 1  \end{aligned} \right\} Q^n $ } & \multirow{2}{*}{\centering 输出状态不变} \\ % 第一行
            0 & 0 & 1 &   &  \\ % 第二行
            \hline
            0 & 1 & 0 & \multirow{2}{*}{ $ \hspace{0.4cm} \left. \begin{aligned} 0 \\ 0 \end{aligned} \right\} 0 $ } & \multirow{2}{*}{\centering 输出状态与$J$相同} \\ % 第三行，合并说明单元格
            0 & 1 & 1 &   & \\
            \hline
            1 & 0 & 0 & \multirow{2}{*}{ $ \hspace{0.4cm} \left. \begin{aligned} 1 \\ 1 \end{aligned} \right\} 1 $ } & \multirow{2}{*}{\centering 输出状态与$J$相同} \\ % 第四行
            1 & 0 & 1 &   & \\
            \hline
            1 & 1 & 0 & \multirow{2}{*}{ $ \hspace{0.7cm} \left. \begin{aligned} 1 \\ 0  \end{aligned} \right\} \overline{Q}^n $ } & \multirow{2}{*}{\centering 每来一个脉冲，输出状态改变一次} \\ % 第五行
            1 & 1 & 1 &   & \\
            \hline
        \end{tabular}
        \caption{主从$J-K$触发器的特性表} % 表格标题
        \label{tab:D_flip_flop} % 表格引用标签
    \end{table}
    % 表格代码结束

\end{document}

